% 鸟神星（综述）
% license CCBYSA3
% type Wiki

本文根据 CC-BY-SA 协议转载翻译自维基百科\href{https://en.wikipedia.org/wiki/Makemake}{相关文章}。

\begin{figure}[ht]
\centering
\includegraphics[width=6cm]{./figures/abca2436367ca8e9.png}
\caption{2015 年 4 月，哈勃空间望远镜拍摄的低分辨率阋神星及其无名卫星 S/2015 (136472) 1 图像} \label{fig_Makema_1}
\end{figure}
玛克马克$^{[g]}$（小行星编号：136472 Makemake）是一颗位于柯伊伯带的矮行星，柯伊伯带是海王星轨道之外由冰质天体组成的圆盘。它是第四大海王星外天体，也是经典柯伊伯带中体积最大的成员$^{[h]}$，其直径约为冥王星的 60\%。它于 2005 年 3 月 31 日由美国天文学家迈克尔·E·“迈克”·布朗（Michael E. "Mike" Brown）、查德·特鲁希略（Chad Trujillo）和大卫·拉比诺维茨（David Rabinowitz）在帕洛马天文台发现。作为该团队发现的最大天体之一，玛克马克的发现促成了 2006 年冥王星被重新归类为矮行星。

玛克马克在表面特征方面与冥王星类似：其表面高度反射，大部分覆盖着冻结的甲烷，并因凝聚物（tholins）而被染成红褐色$^{[i]}$。玛克马克目前已知有一颗卫星，但尚未命名。该卫星的轨道表明玛克马克自转具有很高的轴倾角，这意味着它会经历极端季节。玛克马克显示出地球化学活动与低温火山活动的证据，这使科学家怀疑其可能拥有一个地下液态水海洋。在玛克马克上发现了气态甲烷，但尚不清楚它是存在于大气中，还是来自暂时性的气体逸出。

由于尚未有太空探测器近距离访问玛克马克，因此并不存在关于其表面的高分辨率影像。玛克马克距离地球非常遥远，即使通过望远镜观测，它看起来也只是一颗类似星点的光斑。科学家表达了希望派遣太空探测器访问玛克马克的愿望，因为它具有地质活动及潜在的地下海洋。
\subsection{历史}
\subsubsection{发现}
\begin{figure}[ht]
\centering
\includegraphics[width=8cm]{./figures/7d2c811983023908.png}
\caption{玛克马克是在帕洛马天文台使用 1.22 米塞缪尔·奥辛望远镜（如图）拍摄的图像中被发现的} \label{fig_Makema_2}
\end{figure}
玛克马克于 2005 年由美国天文学家迈克尔·E·“迈克”·布朗（Michael E. "Mike" Brown）、查德·特鲁希略（Chad Trujillo）和大卫·拉比诺维茨（David Rabinowitz）组成的团队发现，他们当时正在寻找海王星轨道之外的大型天体$^{[23]:200}$。该团队对海王星外天体的搜索始于 2001 年$^{[23]:187-188}$，具体做法是使用安装在加利福尼亚州帕洛马天文台塞缪尔·奥辛 1.22 米（48 英寸）望远镜上的 CCD（电荷耦合器件）相机$^{[j]}$，定期对夜空进行成像$^{[25][23]:190}$。玛克马克的发现图像由该望远镜于 2005 年 3 月 31 日拍摄$^{[1][26]}$，但直到 2005 年 4 月 3 日，迈克·布朗在检查这些图像时才发现该天体，并识别出它具有异常高的亮度$^{[27]:132}$。

在发现玛克马克前几个月，布朗及其团队还发现了两个异常巨大的海王星外天体，闇神星（Haumea）与厄里斯（Eris），它们被认为至少与当时的第九行星冥王星体积相当$^{[28]}$。当团队正计划对这两个天体开展进一步观测时，原本计划将玛克马克的公布推迟至 2005 年 10 月厄里斯公布之后$^{[23]:202[27]:133-134}$。然而这一计划被打乱，因为 2005 年 7 月 27 日，由西班牙内华达山天文台 José Luis Ortiz Moreno 领导的团队宣布他们自己的闇神星发现结果$^{[27]:145[23]:207}$。布朗意识到，他所在团队包含闇神星、厄里斯与玛克马克位置的观测日志无意中公开，并被 Ortiz 所在机构的一台计算机访问过$^{[27]:154-155[23]:211}$。担心厄里斯与玛克马克的发现也会被抢先公布，布朗于 2005 年 7 月 29 日联系小行星中心（MPC）的 Brian G. Marsden 宣布发现$^{[29][27]:156}$。MPC 在加利福尼亚当地时间中午在其网站上发布了厄里斯与玛克马克的发现公告，随后当天晚上由中央天文电报局发布消息$^{[23]:210[30][26]}$。这些与冥王星大小相当的天体的公布，引发了关于“什么是行星”的广泛争论$^{[25]}$，这促使国际天文学联合会（IAU）在 2006 年 8 月制定新的行星定义，将冥王星重新归类为矮行星$^{[31][32]}$。
\subsubsection{名称与符号}
这颗矮行星以玛克马克命名，玛克马克是复活节岛本土拉帕努伊人神话中人类创造者与丰饶之神$^{[6]}$。其小行星目录编号为 136472，该编号由小行星中心在 2005 年 9 月 7 日授予，当时其轨道已被充分确定$^{[33][34]}$。在命名之前，玛克马克的临时编号为 2005 FY9，当其发现被公布时由小行星中心赋予$^{[30][6]}$。玛克马克之前还被其发现团队昵称为“Easterbunny”$^{[k]}$，这是引用其发现时间（复活节之后不久），并曾被布朗用于发现记录的自动代码名称“ K05331A ”$^{[23][7]}$。

在其个人文字与访谈中，布朗回忆称，确定玛克马克的命名非常困难，因为当时已知的天体特征与神话无明显联系$^{[7][27]:246[35][36]}$。为了保留该天体与复活节的联系，布朗曾考虑用盎格鲁-撒克逊女神 Ēostre，或阿尼什纳贝族的诡计兔 Manabozho 命名，但发现这两个名称都无法使用$^{[l]}$。最终，布朗与其团队选择了“Makemake”这一名称，满足了天体与复活节的联系，以及 IAU 关于以创世神命名经典柯伊伯带天体的规定$^{[32][37]}$。玛克马克的名称于 2008 年 7 月获 IAU 批准并公布$^{[m]}$。

玛克马克的符号 ⟨🝼⟩ 于 2022 年 1 月作为 U+1F77C 被纳入 Unicode$^{[38]}$。IAU 不鼓励在科学出版物中使用行星符号$^{[39]}$，因此该符号主要由占星学者使用$^{[40]}$。不过 NASA 曾在 2015 年发布的一份信息图中使用过该符号一次$^{[41][40]:4}$。玛克马克符号由 Denis Moskowitz 与 John T. Whelan 设计，它代表玛克马克面部的传统岩刻画，并经过设计使其类似字母“M”$^{[42]}$。其他占星师也设计并使用自己的玛克马克符号，例如 ⟨⟩$^{[40]:5}$。
\subsection{轨道与分类}
\begin{figure}[ht]
\centering
\includegraphics[width=8cm]{./figures/bc08e081cf518854.png}
\caption{示意图显示了玛克马克绕太阳运行的倾斜轨道（灰色），并标示了外行星的位置。沿着玛克马克轨道路径的垂直灰线标出其位于黄道平面之上和之下的位置。} \label{fig_Makema_3}
\end{figure}
玛克马克在海王星之外绕太阳运行，其平均距离为 45.5 个天文单位（AU；68.1 亿公里或 42.3 亿英里）$^{[b]}$。它完成一次公转需要 307 年$^{[b]}$。玛克马克的轨道偏心率为 0.16，遵循一条中等椭圆轨道$^{[44]:212-213}$，距离太阳最近时为 38.2 AU（近日点），最远时为 52.8 AU（远日点）$^{[b]}$。玛克马克相对于黄道的轨道倾角较大，为 29°$^{[9][44]:212}$。

玛克马克目前接近远日点，即其轨道上距离最远的位置$^{[16]:9}$。截至 2025 年 11 月，它距离太阳为 52.7 AU$^{[45][21]}$，并将于 2033 年 5 月抵达远日点$^{[46]}$。玛克马克目前位于黄道平面之上很远的位置$^{[47][44]:212}$，并将在远日点保持这一状态，届时其黄道纬度将达到 25.9°$^{[46]}$。玛克马克将在 2103 年穿过黄道平面$^{[48]}$，并将在 2186 年到达近日点，位于黄道平面以下 26°$^{[11]}$。N 体模拟显示，玛克马克的轨道在数十亿年尺度上是稳定的，在太阳系余下寿命中不太可能发生显著变化$^{[49]:6-7}$。

玛克马克与海王星之外许多小型冰质天体共享其轨道特征，这些天体共同构成了柯伊伯带这一区域。玛克马克特别属于“动力学高温”的经典柯伊伯带天体群$^{[n][44]:212}$，其轨道特征具有较高的倾角（i > 5°）、较低的偏心率（e < 0.2），且与海王星不存在轨道共振$^{[50]:21[51]:2}$。玛克马克是经典柯伊伯带中体积最大的成员$^{[44]:212}$，尽管它仅占该带总质量的一小部分$^{[52]:8[o]}$。动力学高温的经典柯伊伯带天体被认为在太阳系早期历史中被海王星引力散射$^{[50]:22}$，因此天文学家也将玛克马克称为“散射”天体$^{[53][54]:L98[55]:284}$。

科学界的共识是玛克马克是一颗矮行星：即它质量足够大，其自身引力可以使其形成接近球形的形状，但质量又不足以清除轨道上的其他天体，这一点从其位于柯伊伯带即可证明$^{[37][32]}$。玛克马克是 IAU 在新程序下命名的首个预期成为矮行星的天体，也是 2006 年确立该分类以来继原有的谷神星、冥王星与厄里斯之后，第四个被宣布为矮行星的天体$^{[m]}$。玛克马克更具体地说是一颗冥族矮行星（plutoid），即绕海王星之外运行的矮行星子类$^{[6][35]}$。
\subsection{大小、形状与质量}
\begin{figure}[ht]
\centering
\includegraphics[width=8cm]{./figures/c883c7e09b0f9dea.png}
\caption{比较直径大于 700 公里（430 英里）的各大海王星外天体的尺寸、反照率与颜色。玛克马克位于第一行，从右数第二个。深色弧线表示天体尺寸的不确定范围。} \label{fig_Makema_4}
\end{figure}
玛克马克是一个近乎球形的天体，其平均直径约为 1,430 公里（890 英里）$^{[12]:2}$，约为冥王星直径的 60\%（3/5）$^{[p][57]}$，或地球直径的 11\%（1/9）$^{[31]}$。这使得玛克马克成为太阳系中已知的第四大矮行星及海王星外天体，仅次于冥王星、厄里斯和闇神星$^{[58]}$。2011 年对一次恒星掩食的观测显示，玛克马克在极地略呈扁平状，极直径上限约为 1,420 公里（880 英里）$^{[q]}$，赤道直径约为 1,434 公里（891 英里）$^{[12]}$。这些尺寸与一种称为麦克劳林扁球（Maclaurin spheroid）的扁平球形结构相一致，这种形状出现在天体处于流体静力平衡状态（即天体自身引力足以将其压成球形）并因自转而发生形变时$^{[12]:2[17]:5[19]:12}$。

玛克马克的质量估计在约 (2.5\times10^{21}) 至 (2.9\times10^{21}) 千克之间，该数值由其卫星的轨道周期与距离确定$^{[15]:3}$。这使得玛克马克成为太阳系中已知质量第四大的矮行星及海王星外天体，仅次于厄里斯、冥王星和闇神星$^{[59]}$。与太阳系其他天体相比，玛克马克的质量约为地球月球质量的 3.7\%（或地球质量的 0.045\%）$^{[r]}$，约为冥王星质量的 20\%$^{[s]}$。根据玛克马克的质量和平均直径，其平均表面重力约为 0.35 m/s(^2)$^{[e]}$（约为地球重力的 3.6\%）$^{[t]}$，其表面逃逸速度约为 0.71 km/s$^{[f][16]:8[ ]}$。

\subsection{自转}
玛克马克的自转周期尚不确定，截至 2025 年的测量显示，该周期为 11.4 或 22.8 小时（0.48 或 0.95 日）$^{[19]:2,7}$。这些自转周期的测量通过监测玛克马克亮度随时间的变化进行，并绘制其光变曲线$^{[17][19]:2}$。玛克马克的亮度变化非常微小（0.03 星等），推测是由于其表面反照率变化较小，这使得望远镜很难测量玛克马克的光变曲线与自转周期$^{[17]:1,6}$。例如，2019 年之前的研究提出了可能的自转周期 7.77、11.24、11.5 与 22.48 小时$^{[17]:1}$。截至 2025 年的测量尚不清楚玛克马克的亮度在一次自转中是出现一次还是两次峰值，因此不确定其自转周期究竟是 11.4 小时，还是其两倍即 22.8 小时$^{[19]:2}$。

玛克马克的自转轴倾角尚未被测量，不过可以合理推测其自转轴与其卫星轨道的极点对齐$^{[8]:4[61]:16}$。在这种情况下，玛克马克的自转轴相对于绕太阳轨道的倾角介于 46° 至 78° 之间（或相对于黄道平面介于 63°–87° 之间），在其卫星被发现时，其赤道朝向太阳与地球（接近春分）$^{[8]:3-4}$。如此高的轴倾角加上其偏心轨道，可导致玛克马克表面温度与地形出现显著季节变化，类似于冥王星$^{[8]:4-5[61]:16}$。玛克马克的卫星被预测将在 2009–2013 年或 2023–2027 年期间掩食玛克马克，因此如果其自转与其卫星轨道一致，玛克马克可能已经在这两个年份范围之一经历过春分$^{[15]:1}$。
\subsection{地质}
\subsubsection{表面}
\begin{figure}[ht]
\centering
\includegraphics[width=10cm]{./figures/a8a10a028dc3e3f6.png}
\caption{詹姆斯·韦布空间望远镜测得的玛克马克近红外光谱。甲烷（$CH_4$）在 1–2 微米范围内的吸收特征在玛克马克的光谱中极为显著，说明其表面甲烷含量非常丰富。玛克马克上检测到的其他化合物包括乙烷（$C_{2}H_6$）、乙炔（$C_2H_2$）、氘代甲烷（$CH_3D$），以及可能存在的乙烯（$C_2H_4$）$^{[16]:2}$。} \label{fig_Makema_5}
\end{figure}
\begin{figure}[ht]
\centering
\includegraphics[width=8cm]{./figures/6453132cf5651cb1.png}
\caption{玛克马克的艺术示意图，呈现其均匀的浅棕色表面以及缺乏显著大气的特征} \label{fig_Makema_7}
\end{figure}
由于距太阳非常遥远，玛克马克表面的极低温度约为 30 至 40 K（−243 至 −233 ℃；−406 至 −388 ℉）$^{[19]:5[16]:17}$——冷到足以让某些挥发性物质（如甲烷）以固态冰形式存在$^{[62]}$。天文光谱研究显示，玛克马克表面主要由冻结的甲烷构成，同时含有少量长链烃类，包括乙烷$C_2H_6$）、乙烯（$C_2H_4$）、乙炔（$C_2H_2$）以及多种高分子量烷烃（如丙烷）$^{[63][16]:2}$。

在可见光范围内，玛克马克表面显得非常明亮且具有很高反射率，其几何反照率为 82\%（比冥王星更高）$^{[17]:7[18]:5}$，这表明其甲烷可能是新近沉积的$^{[61]:15[64]:3-4}$。玛克马克的甲烷冰在近红外波段具有很强的吸收性，说明它可能以异常大的厘米级颗粒存在，或者更可能是由烧结颗粒构成的厚板状结构$^{[31][65]:3597[64]:16}$。与此同时，“新视野号”探测器的相位曲线测量指出，玛克马克表面的风化层由光滑的颗粒组成，类似于雪$^{[18]:20}$。

玛克马克表面的长链烃来源于甲烷受到紫外阳光及宇宙射线照射，其过程中甲烷被分解并触发光化学反应$^{[63]:1[16]:2,9}$。这些光化学反应可以级联进行：将甲烷转化为乙烷，再转化为乙烯、乙炔，依次类推$^{[65]:3594-3595}$，最终生成一种暗色、红褐色的复杂烃类混合物，称为凝聚物（tholins）$^{[55]:285[64]:4}$。这些凝聚物赋予玛克马克红褐色的外观$^{[i]}$，与冥王星表面观察到的现象类似$^{[31][67]}$。玛克马克的颜色比冥王星略浅，但比厄里斯略红$^{[18]:5}$；这种颜色差异可能源于不同矮行星表面凝聚物含量不同$^{[20]:5475}$。尽管凝聚物应当使玛克马克表面变暗，但由于新鲜甲烷冰覆盖了这些凝聚物，该矮行星仍然保持明亮$^{[22]:569[61]:15[64]:3-4}$。

玛克马克与冥王星和厄里斯一样拥有高含量的甲烷冰，但与后两者不同的是，玛克马克表面似乎缺乏一氧化碳和氮冰$^{[64]:1-2}$。詹姆斯·韦布空间望远镜（JWST）未能在玛克马克表面发现这两种冰，说明其中氮含量低于 3\%，一氧化碳含量低于百万分之一$^{[64]:13[16]:1}$。由于缺乏可与之混合的氮和一氧化碳，玛克马克上的甲烷冰保持纯净，并可形成很大的厚度或颗粒尺寸$^{[55]:288[64]:16}$。玛克马克缺乏氮是预期之中，因为氮高度挥发，其蒸气较之从冥王星和厄里斯较强引力逃逸，更易从玛克马克的引力逃逸$^{[62]:287[61]:16[64]:18}$。玛克马克明显缺乏一氧化碳的原因则不太明确：它可能是通过大气逃逸或由内部热液驱动的地球化学反应而被移除，或玛克马克在形成时就含有较低含量的一氧化碳$^{[64]:19[70]:10[71]:5}$。水冰和二氧化碳冰在玛克马克表面似乎也不存在，尽管它们是柯伊伯带天体中常见的不挥发性材料；这可能是因为在玛克马克上，这些冰完全被甲烷及其辐照产物覆盖$^{[65]:3598[70]:13}$。

玛克马克表面似乎非常均一，在反照率、颜色和成分上经度差异极小$^{[72][65][17]:1,6}$，这与冥王星高度斑驳的地形形成鲜明对比$^{[61]:16}$。尚不清楚玛克马克是否存在纬向表面差异，因为检测这类差异需要在玛克马克绕太阳运行（换言之，出现季节变化）时，长时间连续观测玛克马克视角角度的改变$^{[u]}$，而这需要多年$^{[17]:6-7[19]:8}$。从 2006 到 2017 年间，玛克马克的绝对星等与光变曲线未发生变化，而此期间玛克马克视角角度改变了约 11°$^{[17]:7}$。如果玛克马克存在纬向表面差异，它们可能呈现沿经度方向分布的条带状结构$^{[61]:16}$。行星科学家 William M. Grundy、Alex H. Parker 及其同事推测，玛克马克丰富的挥发性甲烷可能导致与冥王星类似的地貌与地质特征$^{[61]:16[64]:4}$。如果玛克马克存在与冥王星类似的季节性挥发物迁移过程，它可能形成一条沿经度方向均一的暗色物质带，类似冥王星的 Belton 区域$^{[61]:16}$。或者，如果玛克马克曾拥有非全球性大气并冻结在表面，其赤道可能呈现明亮、霜状覆盖，而两极则较暗$^{[61]:16}$。季节性的甲烷升华与沉积可能产生刀刃状地形，甚至厚重、类似冥王星斯普特尼克平原（Sputnik Planitia）的对流冰川$^{[64]:4}$。玛克马克不太可能拥有高于 10 公里（6.2 英里）的山脉$^{[17]:6}$。
\subsubsection{内部结构与可能的地质活动}
玛克马克的整体密度约为 $1.67 g/cm^3$（不确定度为 $\pm0.17 g/cm^3$）$^{[15]:3}$，与海王星外矮行星冥王星、公公（Gonggong）和夸欧尔（Quaoar）相似$^{[59]:7}$。与这些矮行星一样，该密度表明玛克马克的内部主要由水冰和岩石构成$^{[59]:7[73]:10}$。玛克马克体积足够大，其内部很可能是分化的，具有一个岩石核心，并由冰层包围$^{[74]:230[75][73]:8}$。行星科学家推测，玛克马克内部可能含有足够的放射性核素和原初热量，可以在过去，甚至可能在如今，维持一个地下液态水海洋$^{[73][76]}$。大量内部热量可能引发地质现象，例如低温火山活动$^{[77][19]:4}$。

JWST 的光谱观测在玛克马克表面检测到含有氘（D 或 ${^2}H$）和碳-13（${^13}C$）的甲烷重同位素，其中，天文学家测得玛克马克的氘氢比（D/H）为 ($(2.9\pm0.6)\times10^{-4}$)，碳-13 与碳-12（$^{13}C/^{12}C$）比例为$0.010\pm0.003$.$^{[64]}$。虽然玛克马克的 $^{13}C/^{12}C$比例与其他太阳系天体一致，但 D/H 比例有所不同：它远低于彗星中甲烷的 D/H 值，但与彗星中水的 D/H 值相似$^{[64][78]}$。行星科学家将玛克马克较低的 D/H 比例解释为其内部具有温暖环境并存在活跃热液地球化学活动的证据：玛克马克贫氘甲烷可能从地下水的地球化学反应中获得氢元素，而这些反应需要约 150 ℃（302 ℉）的高温，只能由玛克马克假定核心的热量维持$^{[78][73]}$。在该情景下，玛克马克的地下水可能以液态水或对流固态冰的形式存在，而内部产生的甲烷可能通过逸气作用或低温火山喷发运输到玛克马克表面$^{[73]}$。然而，玛克马克贫氘甲烷也可能是原初的（直接来源于原始太阳星云的吸积），因此内部地球化学活动并非解释其存在的必要条件$^{[71]}$。

与远红外相比，玛克马克发出大量中红外辐射，这一现象自 2008 年由斯皮策空间望远镜首次报告以来，引发了天文学界的多种解释$^{[19]:1}$。天文学家最初认为，玛克马克过量的中红外辐射来自混合在明亮、寒冷地形中的暗色、温暖地形（以及在其卫星被发现之后，也可能来自卫星），但该假说无法准确描述玛克马克在不同波长下的红外辐射$^{[19]:1-2}$，也无法解释玛克马克极小的亮度变化$^{[79]}$。2025 年，Csaba Kiss 及其合作者提出，玛克马克过量中红外辐射可能由温度约 150 K（−123 ℃；−190 ℉）的低温火山热点引起，或者由环绕其轨道运行的微小碳质尘埃环造成$^{[19]:1-2[76]}$。低温火山热点情景更为合理，因为上述尘埃环很容易因太阳辐射压力而失稳，尽管如果低温火山喷发能够将碳质尘埃抛射到玛克马克轨道中，该尘埃环仍可能持续补给$^{[19]:6}$。所提出的低温火山热点可能释放的热量与土星卫星土卫二（Enceladus）南极间歇泉相当，并可能喷发含有氨与各种溶解盐的液态水低温熔岩$^{[19]:4[76]}$。该低温火山热点在玛克马克表面的位置未知，但估计其覆盖面积约为 $350 km^2$（140 平方英里，约等同于半径 ~10 km 或 6.2 英里的圆）$^{[19]:3-4}$。
\subsection{大气或逸气}
\begin{figure}[ht]
\centering
\includegraphics[width=8cm]{./figures/9b3ce67a8a26923c.png}
\caption{JWST 在玛克马克的近红外光谱中检测到气态甲烷（$CH_4$）的荧光（左图，标记为 a）。观测到的荧光可以用甲烷逸出所形成的彗发（b），或一层稀薄的甲烷大气（c）来解释。} \label{fig_Makema_8}
\end{figure}
2025 年对 JWST 光谱的分析揭示，玛克马克上存在气态甲烷，其在近红外波段因吸收阳光而产生荧光$^{[80][16]}$。玛克马克是在冥王星之后第二个被确认存在气体的海王星外天体$^{[16][80]}$。然而，目前尚不清楚玛克马克的甲烷气体是被包含在一个由引力束缚的大气层中，还是由于甲烷冰升华或低温火山羽流从其表面暂时逸出（甚至逃逸）$^{[16]:1[80]}$。理论上，玛克马克的质量和温度仅勉强足以维持由甲烷或氮构成的大气；JWST 的观测显示玛克马克似乎不含氮气，这表明大部分氮已经通过大气逃逸而损失$^{[16]:17}$。

如果 JWST 检测到的甲烷气体全部包含在一个由引力束缚的大气中，那么其表面大气压约为 10 皮巴（1 微帕），比地球大气压低 1000 亿倍，比冥王星大气压低 100 万倍$^{[80]}$。如此极其稀薄的大气并未在 2011 年玛克马克恒星掩食观测中被探测到，这也支持掩食观测的结论，即玛克马克缺乏大于 4–12 纳巴（0.4–1.2 毫帕）级别的全球性大气$^{[16]:9[22][65]}$。这种假定稀薄大气的温度约为 40 K（−233.2 ℃；−387.7 ℉），略高于该大气压下甲烷的升华温度。这意味着玛克马克的假定大气可能由表面甲烷冰的升华维持$^{[80][16]:9}$。由于玛克马克具有偏心轨道，其假定大气可能随与太阳距离的变化而改变：例如，在近日点较高的温度下，玛克马克可能升华更多甲烷，但也可能损失部分甲烷至大气逃逸$^{[62]:L62}$。

另一种可能是，JWST 检测到的甲烷气体仅来自逸气，那么这将意味着玛克马克每秒从其 4–30\% 的总表面积释放约 266 千克（586 磅）甲烷$^{[16]:9}$。目前尚不清楚这些甲烷逸气的速度是否足以逃离玛克马克的引力。如果甲烷气体正在逃逸，它会形成类似彗星的彗发包围玛克马克$^{[16]:8}$。估算的质量损失速率可与土卫二（Enceladus）的水羽流（300 千克/秒或 660 磅/秒）相当，而甲烷释放的有限表面积可能与先前提出的玛克马克低温火山热点有关$^{[16]:9}$。研究人员推测，甲烷的低温火山逸气在玛克马克等海王星外矮行星中可能普遍存在$^{[81]:5}$。
\subsection{卫星与潜在光环}
\subsubsection{S/2015 (136472) 1}
\begin{figure}[ht]
\centering
\includegraphics[width=8cm]{./figures/a55684c7c37878b5.png}
\caption{2015 年 4 月，哈勃空间望远镜拍摄的玛克马克卫星发现图像。该卫星在 4 月 27 日可见，但到 4 月 29 日已移动并被隐藏。} \label{fig_Makema_9}
\end{figure}
玛克马克目前已知只有一颗天然卫星，尚未命名，临时编号为 S/2015 (136472) 1，非正式昵称为“MK2”$^{[8][75]}$。它由天文学家 Alex H. Parker、Marc W. Buie、William M. Grundy 与 Keith S. Noll 在 2015 年 4 月 27 日哈勃空间望远镜拍摄的图像中发现，并于 2016 年 4 月 26 日宣布$^{[47]}$。在可见光中，S/2015 (136472) 1 比玛克马克暗约 1,300 倍（7.8 星等），并推测其表面非常暗，直径约为 175 公里（109 英里），以解释玛克马克过量的中红外辐射$^{[79][8]:3-4}$。该卫星很可能沿近圆轨道绕玛克马克运行，其轨道周期为 18 天，半长轴为 $22,250 \pm 780$ 公里（$(13,830 \pm 480)$ 英里）$^{[15]}$。

当 S/2015 (136472) 1 被发现时，其轨道从地基观测点视角几乎呈边缘朝向，这意味着该卫星在视野中似乎穿过玛克马克前方或背后$^{[8]:2[15]:3}$。尽管这种边缘构型使望远镜难以成像 S/2015 (136472) 1$^{[79]}$，但可能使其能够发生掩食与遮蔽玛克马克的现象$^{[8]:4[75]}$。预测该卫星可能在 2009–2013 年期间掩食过玛克马克，或可能在 2023–2027 年期间仍在掩食玛克马克$^{[15]:1}$。截至 2025 年，尚未有关于 S/2015 (136472) 1 掩食事件的报告$^{[15]}$。
\begin{figure}[ht]
\centering
\includegraphics[width=10cm]{./figures/dcd76d3668b9b3dc.png}
\caption{} \label{fig_Makema_10}
\end{figure}
\subsubsection{可能存在其他卫星}
哈勃空间望远镜的成像观测显示，在距离玛克马克 30,000 公里（19,000 英里）以外，它没有比视星等 26.9（约比玛克马克暗 10 等）$^{[v]}$更亮的额外卫星 $^{[82]:8}$。如果有较大的卫星绕玛克马克非常靠近运行，它们可能会被望远镜的视场所隐藏$^{[17]:6}$。天文学家曾讨论过玛克马克存在一个比 S/2015 (136472) 1 更大的暗卫星的可能性，以解释玛克马克过量的中红外辐射与其明显较慢的自转$^{[17]:6}$，但这一假设并不受支持，因为这需要一个规模不现实地巨大的卫星$^{[19]:3}$。
\subsubsection{可能存在环系统}
玛克马克目前没有已知的光环。对于遥远天体而言，光环过小、过暗，无法直接被望远镜成像，因此理想的探测方式是利用恒星掩食观测$^{[83]:27}$。然而，在 2011 年玛克马克发生的恒星掩食观测中，并未发现光环。如果玛克马克确实存在光环，它们很可能沿其赤道以类似于 S/2015 (136472) 1 的边缘朝向构型运行，这可能使它们在 2011 年的掩食事件中未被天文学家探测到$^{[19]:5}$。

玛克马克可能存在光环的假说，曾被用作解释其过量中红外辐射的潜在方案，但被认为不太可能，因为这一假设要求光环由极小（约 100 纳米）尘埃颗粒构成，而这类尘埃会在十年内被太阳辐射压破坏$^{[19]:6,15}$。尽管如此，如果玛克马克存在“牧羊卫星”、由碰撞碎片及小卫星连续产生尘埃，或者存在将尘埃喷射入轨道的低温火山喷发，仍有可能维持这样的光环$^{[19]:6}$。
\subsection{起源}
与柯伊伯带中的其他矮行星类似，玛克马克被认为形成于约 45 亿年前的早期太阳系历史$^{[31][84]}$。柯伊伯带矮行星被假设最初是小型微行星，在几百万年间通过吸积周围物质及其他微行星而增长到目前的大小$^{[74]:214}$。玛克马克形成时的环境温度必须足够低，才能使甲烷等挥发物凝结成固体并进一步吸积到矮行星之中$^{[64]:3,17[71]:3}$。不过，玛克马克在吸积阶段可能失去了一部分原始甲烷，因为它最初质量较小，并由于频繁撞击事件和更强的太阳辐射而温度更高$^{[73]:9}$。此外，还有假说认为，在玛克马克的过去某个时期，与另一天体的大规模碰撞可能形成了其卫星 S/2015 (136472) 1$^{[75]}$。

根据 2020 年基于太阳系形成模型（对 2005 年 Nice 模型的更新、后者首次提出类似情景）提出的一项假说，在太阳系形成数千万年后，巨行星之间的引力相互作用使海王星突然向外迁移至太阳 15–30 AU 之间的一片巨大环星盘中，从而引力散射了其中的大量天体$^{[85][86]:176}$。该模型显示，包括玛克马克在内的几乎所有柯伊伯带天体最初都形成于比现在更靠近太阳的那片环星盘中$^{[86]:175-76[73]:9}$。这一盘的散射被认为形成了现今柯伊伯带中的共振族与“动力学热”$^{[n]}$古典族（玛克马克即位于此），以及散射盘$^{[86]:176}$。
\subsection{观测与探索}
\subsubsection{观测}
\begin{figure}[ht]
\centering
\includegraphics[width=6cm]{./figures/ca03c482bf875fd0.png}
\caption{2004年12月13日斯隆数字巡天（SDSS）对马凯马凯的追溯发现图像（圆圈标出），其上方为大型星系NGC 4274。} \label{fig_Makema_11}
\end{figure}
从视觉绝对星等来看，Makemake 是仅次于厄里斯和冥王星的第三亮已知海王星外天体$^{[87]}$。Makemake 的高度本征亮度源于其体积较大以及表面具有极高反射率$^{[w]}$。另一方面，从视觉视星等来看，Makemake 是天空中仅次于冥王星的第二亮海王星外天体，这是因为它与太阳的距离比厄里斯更近$^{[25][68]}$。每年 3 月至 4 月进入冲日位置时，Makemake 的视星等峰值约为 17 等$^{[89][90]}$，在这种亮度下，可使用高级业余望远镜观测到它$^{[6]}$。由于 Makemake 与地球距离极为遥远，在望远镜中它呈现为一个角直径约 38 毫角秒的微小天体$^{[22]:568}$，因此望远镜无法将其分辨成除类似恒星光点之外的形状$^{[89]}$。在夜空中，Makemake 位于北天的后发座，自其被发现以来一直位于该星座$^{[90]}$。在 2028 年末，Makemake 将移动至牧夫座$^{[90]}$。

尽管是最明亮的海王星外天体之一，Makemake 却被发现得相对较晚 — 远远晚于许多更暗的海王星外天体$^{[44]:212}$。这是因为 Makemake 具有高度倾斜的轨道，使其远离黄道平面，而此前的巡天主要集中在黄道附近$^{[44]:212[91]:1}$。虽然若干巡天项目在 Makemake 被正式发现前的数年内曾意外拍摄到它（这类观测称为 “预发现”），但这些影像直到事后才被识别出来$^{[1]}$ $^{[x]}$。目前已知最早的 Makemake 预发现影像来自 1955 年 1 月 29 日帕洛马山天文台拍摄的一张天文底片，该影像比 Makemake 的正式发现早了 50 多年（约占其轨道周期的 16%）$^{[1]}$。

在天空中移动时，Makemake 可能从地球视角掩食一颗背景恒星，短暂遮挡其光线，从而产生恒星掩星现象$^{[96]}$。Makemake 的恒星掩星有助于揭示其形状及潜在大气特征，但由于该矮行星距离地球极为遥远，其位置存在较大不确定性，因此极难准确预测$^{[97][22]:566}$。Makemake 所在天空区域恒星密度较低，使其恒星掩星事件极为罕见$^{[96]}$。截至 2025 年，仅有一次 Makemake 的恒星掩星被成功预测并观测到$^{[98]:5}$。首次也是唯一一次观测到的 Makemake 恒星掩星发生在 2011 年 4 月 23 日，当时位于南美各地的 16 台望远镜中有 7 台获得了有效观测数据$^{[96][22]:566}$。
\subsubsection{探测}
\begin{figure}[ht]
\centering
\includegraphics[width=6cm]{./figures/2fb1af0b7c4aa39f.png}
\caption{Makemake（以红色标记显示）由“新视野”号探测器于 2007 年 10 月 6 日拍摄。} \label{fig_Makema_12}
\end{figure}
玛克马克尚未被太空探测器近距离造访，尽管天文学家与行星科学家已经表达了向其发射探测器的愿望$^{[73]}$。玛克马克被视为一个具有吸引力的探索目标，因为它可能存在具有持续地质活动的地下海洋$^{[99][73:13]}$。对像玛克马克这样的海王星外天体进行探索，有助于深入理解太阳系的形成与演化$^{[100][101]}$。

2011 年，Ryan McGranaghan 及其同事的一项研究计算出：基于 2036 年 8 月 24 日的发射日期，利用木星引力助推，对玛克马克的飞掠任务耗时仅略超过 16 年。探测器抵达时，玛克马克距离太阳约为 52.3 AU$^{[100:300]}$。2024 年，美国田纳西大学的一项研究指出，如果飞掠任务采用有动力的木星引力助推，则可在更短的 9.6–16.4 年内抵达玛克马克，具体时间取决于航天器的载荷质量$^{[101:17]}$。对于 2036 年 8 月 22 日和 2048 年 9 月 27 日的发射日期，有动力的木星引力助推将是最优方案$^{[101:7]}$。

2019 年，Amanda Zangari 及其同事的一项研究为不同的引力助推与发射能量设计了多种可能的玛克马克飞掠轨道。若在 2025–2027 年或 2036–2039 年发射，一次木星引力助推分别可在 12.8–23.6 年或 11.6–19.2 年内将航天器送抵玛克马克$^{[102:922]}$。一次土星引力助推可为低能量发射提供更快捷的路线：若在 2032–2033 年或 2036–2040 年发射，航天器分别可在 19.2–22.5 年或 12.8–19.1 年内到达玛克马克$^{[102:923,925]}$。若在 2037–2049 年发射，利用木星与土星的联合引力助推，航天器可在 16.8–17.3 年内抵达玛克马克$^{[102:923]}$。

2007 年 10 月与 2017 年 1 月，“新视野”号探测器分别在 52 AU 与 70 AU 的距离处对玛克马克进行了远距离观测$^{[18:6]}$。探测器穿越柯伊伯带的外向轨道，使得能够在从地球无法达到的高相位角条件下观测玛克马克，从而得以确定玛克马克表面的光散射特性与相位曲线行为$^{[18]}$。
\subsection{参见}
\begin{itemize}
\item 天文命名惯例
\item 太阳系天体尺寸列表
\item 《我如何“杀死”冥王星，以及它为何活该如此》
\item 海洋世界
\item 地外液态水
\end{itemize}
\subsection{注释}\\
a.拉帕努伊语发音为 [ˈmakeˈmake]，在英式英语中英文化读法为 /ˈmækiˈmæki/$^{[2]}$，在美式英语中为 /ˈmɑːkiˈmɑːki/ 或 /ˈmɑːkeɪˈmɑːkeɪ/$^{[3]}$ $^{[4]}$。/ɑː/ 与 /æ/ 的区别反映了英式与美式英语对于波利尼西亚语元音 “a” 的不同处理方式（类似意大利语 “pasta” 中第一个 “a” 的发音）；/eɪ/ 的读法尝试逼近波利尼西亚语元音 “e”$^{[5]}$。\\
b.这些轨道要素以太阳系质心（SSB）为参考系来表达$^{[9]}$。由于行星摄动，太阳围绕太阳系质心在不可忽略的距离上运行，因此日心坐标系下的轨道要素与距离（例如 JPL 小天体数据库中给出的数值$^{[10]}$）会在短时间尺度上发生变化$^{[43]}$。\\
c.使用 $\frac{(a - b)}{a}$ 及 Brown（2013）中的尺寸数据计算得出$^{[12]}$。\\
d.使用 Brown（2013）中的尺寸数据$^{[12]}$，并假设为一扁球体进行计算。\\
e.表面重力（单位为米每二次方秒 m/s²）依据
$\frac{GM}{r^{2}}$
进行计算，其中$G = 6.6743\times10^{-11} m^3\cdot kg^{-1}\cdot s^{-2}$.$^{[60]}$ 为引力常数，$M$ 为玛克马克的质量（单位为千克），$r$ 为玛克马克的半径（单位为米）。\\
f.表面逃逸速度（单位为米每秒 m/s）依据
$v_{esc} = \sqrt{\frac{2GM}{r}}$
进行计算，其中 $G = 6.6743\times10^{-11}m^3\cdot kg^{-1}\cdot s^{-2}$.$^{[60]}$ 为引力常数，$M$ 为玛克马克的质量（单位为千克），$r$ 为玛克马克的半径（单位为米）。\\
g.发音分为四个音节，重音落在两个 a 上。元音的读法因情况而异，参见信息栏。
h.经典柯伊伯带不包括与海王星存在轨道共振的海王星外天体，如冥王星和妊神星。如果将共振天体计算在内，冥王星将是柯伊伯带中最大的天体。\\
i.在行星天文学中，“红色”一词用于描述在较长（更红）波长处反射更多光的天体$^{[66:145]}$。天文学家描述玛克马克与冥王星同样“红”$^{[67:L38]}$ $^{[68:437]}$，尽管对人眼而言，冥王星呈棕色$^{[69]}$。\\
j.发现玛克马克的电荷耦合器件相机为类星体赤道巡天（QUEST）相机，其视场约为 8.3 平方度，具有 1.61 亿像素的分辨率$^{[24]}$。它于 2003 年安装在帕洛玛天文台的塞缪尔·奥申望远镜上，当时 Brown 及其团队对海王星外天体的搜索仍在持续$^{[25]}$ $^{[23:191]}$。\\
k.科学作者 Govert Schilling 报道，Brown 最初曾开玩笑想给玛克马克起绰号“Dead Pope”（死教皇），以影射当时即将去世的约翰·保罗二世教皇，但在其妻子 Diane 的劝阻下，他选择了更不具争议的昵称 “Easterbunny”（复活节兔）$^{[23:202]}$。\\
l.Ēostre 这个名字已用于小行星 343 Ostara，因此根据国际天文学联合会关于名称不得重复的规定，不能再次使用。而对于 Manabozho 这一名称，Brown 个人认为不喜欢，因为其词尾含有 -bozo$^{[32]}$ $^{[7]}$ $^{[27:246]}$。\\
m.国际天文学联合会发布宣布玛克马克命名及其被归类为矮行星的新闻稿于 2008 年 7 月 19 日公布$^{[6]}$，但包括《新科学家》、Space Daily 与行星学会在内的其他来源则在几天前已有报道$^{[35]}$ $^{[32]}$ $^{[5]}$。\\
n.此处的“热”并不指温度（玛克马克非常寒冷），而是指其轨道动力学特性，即轨道受到强烈扰动。\\
o.Jean-Marc Petit 等（2023）估计，热经典柯伊伯带的总质量为 0.012 个地球质量（$\approx7.17\times10^{22}kg$）$^{[52:8]}$。鉴于玛克马克的质量为 $2.69\times 10^{21}kg$.$^{[15]}$，因此认为它约占热经典柯伊伯带总质量的 3.8\%。\\
p.冥王星的直径为 2376 公里$^{[56]}$。\\
q.由于玛克马克两极相对于地球视线的确切扁率与倾角未知，因此 2011 年掩星观测所得的 $1420(_{+18}^{-24})$ 公里这一视极直径仅代表玛克马克真实极直径的上限。这是因为在恒星掩食观测中，只能看到掩食天体的阴影$^{[12]}$。\\
r.地球的质量为 $5.9722\times10^{24}kg$，其卫星（月球）的质量为 $7.346\times10^{22}kg$。\\
s.冥王星的质量为 $1.303\times10^{22}kg$,$^{[56]}$。\\
t.地球的标准表面重力约为 $g_0 = 9.81m/s^2$。\\
u.太阳系天体的视角（aspect angle）定义为该天体自转轴与地球视线之间的夹角$^{[103]}$。\\
v.已知亮度的天体，其大小取决于其反照率。如果一颗比玛克马克暗 10 个星等的卫星，其反照率为 0.7，则其直径大约为 16 公里（9.9 英里）$^{[82:8]}$。反照率越低，其直径则越大。\\
w.太阳系天体的绝对星等（H）根据如下方程计算：
$$H = -5 \log_{10}!\left(\frac{D \sqrt{p}}{1329,\mathrm{km}}\right)~$$，
该式由 Harris & Harris（1997）给出的原始绝对星等与反照率（p）求直径（D）的公式
$$D = \frac{1329,\mathrm{km}}{\sqrt{p}} \times 10^{-0.2H}~$$
变形而来$^{[88:451]}$。在绝对星等方程中，直径 (D) 与反照率 (p) 的数值越大，得到的 (H) 值越小（即越亮）。Mike Brown 在其描述厄里斯发现过程的网页中提到这一关系，他指出矮行星若想显得明亮，可以因为它体积大、反照率高，或两者兼具$^{[25]}$。玛克马克既巨大又高度反射$^{[17]}$。\\
x.Kleť 天文台、帕洛玛天文台与哈雷阿卡拉天文台的追溯观测（precovery）在 2005 年玛克马克宣布后几个月内报告给小行星中心$^{[92][93][94]}$，而斯隆数字巡天（SDSS）的追溯观测则在 2015 年更晚的时候才报告$^{[95]}$。
\subsection{参考资料}
\begin{enumerate}
\item "(136472) Makemake = 2005 FY9"。小行星中心。已存档从原来的2025年10月21日。已检索10月21日，2025。
\item “Makemake” 的定义"。柯林斯英语词典。哈珀柯林斯。2020年2月13日。已存档从原来的2025年11月2日。已检索11月2日，2025。
\item 布朗，迈克 (2008)。"制作-制作"。迈克·布朗的行星。存档自原来的2008年7月17日。已检索7月14日，2008。
\item "Makemake"。Merriam-Webster.com词典。梅里亚姆-韦伯斯特。
\item 艾米丽·拉克达瓦拉(2008年7月15日)。“欢迎来到太阳系，Makemake”。planetary.org。行星社会。已存档从原来的2020年9月24日。已检索11月2日，2025。
\item 克里斯滕森，拉尔斯·林德伯格 (2008年7月19日)。第四颗矮行星叫Makemake(新闻稿)。国际天文学联合会。存档自原来的2008年7月23日。已检索7月20日，2008。
\item 布朗，迈克 (2008)。名字是什么？(第二部分)。迈克·布朗的行星。存档自原来的2020年5月13日。已检索7月14日，2008。
\item 帕克，A.H、; Buie，M。W.;格伦迪，W。M、；诺尔，K。S。(2016年4月25日)。“发现Makemakean月亮”。天体物理学杂志。825(1): l9。arXiv:1604.07461。Bibcode:2016ApJ...825L...9P。doi:10.3847/2041-8205/825/1/L9。S2CID119270442。 
\item “JPL地平线136472 Makemake (2005 FY9) 的在线星历，时代JD 2461000.5 (2025-Nov-21)”。JPL Horizons在线星历系统。喷气推进实验室。已检索10月18日，2025。使用太阳能系统的解决方案重心。星历类型: 元素和中心: @ 0)
\item “JPL小车身数据库查找: 136472 Makemake (2005 FY9)”(2025-08-09最后一次检查)。喷气推进实验室。已检索10月21日，2025。
\item “JPL Horizons 136472 Makemake (2005 FY9) 的在线星历，从2186年1月1日到2187年12月1日”。JPL Horizons在线星历系统。喷气推进实验室。已存档从原始的2021年9月25日。已检索11月1日，2025。(当径向速度或 “rdot” 从负变为正时，就会发生近日点。“hecl-lat” 是日心黄道纬度,或黄道上方的角距离。)
\item M.E.布朗 (2013)。“关于矮行星Makemake的大小、形状和密度”。天体物理学杂志快报。767(1): L7(5pp)。arXiv:1304.1041v1。Bibcode:2013年apj... 767L...7B。doi:10.1088/2041-8205/767/1/L7。S2CID12937717。 
\item "表面椭球体717x717x710-wolfram-alpha"。已存档从原来的2019年12月22日。已检索12月22日，2019。
\item "体积椭球体717x717x710-wolfram-alpha"。已存档从原来的2019年12月22日。已检索12月22日，2019。
\item 丹尼尔·班贝格 (2025)。“矮行星 (136472) Makemake卫星的初步轨道”。arXiv:2509.05880[astro-ph.EP]。
\item 普罗托帕帕，西尔维亚; 黄，伊恩; 莱卢希，伊曼纽尔; 约翰逊，佩里安妮·E；威廉·M·格伦迪；克里斯托弗·R·格莱因等人 (2025年9月)。“JWST在Makemake上检测碳氢化合物和甲烷气体”。天体物理学杂志快报。991(2): l34。arXiv:2509.06772。Bibcode:2025ApJ...991L..34P。doi:10.3847/2041-8213/adfe63。
\item T.A.Hromakina; I.N.Belskaya; Yu。N.克鲁格利; V.G.舍甫琴科; J.L.奥尔蒂斯; P.桑托斯-桑兹; R。Duffard; N.莫拉莱斯; A.Thirouin; R.是的.Inasaridze; V.R。艾瓦齐安; V.T.朱朱纳泽; D.Perna; V.V。鲁缅采夫; 我。V。Reva; A.V。塞雷布里安斯基; A.V。谢尔盖耶夫; 我。E.莫洛托夫; V.A.Voropaev; S.F.Velichko (2019年4月9日)。“矮行星 (136472) Makemake的长期光度监测”。天文学与天体物理学。625: a46。arXiv:1904.03679。Bibcode:2019A & A...625A..46H。doi:10.1051/0004-6361/201935274。S2CID102350991。 
\item 安妮·J·维比瑟；保罗·赫尔芬斯坦; 西蒙·B·波特；苏珊·D·贝内基；卡夫拉斯，J。J.;劳尔，Tod R.; 等人 (2022年4月)。“从新视野观测到的矮行星和大型KBO相位曲线的不同形状”。行星科学杂志。3(4): 95。Bibcode:2022PSJ.....3...95V。doi:10.3847/PSJ/ac63a6。S2CID248448259。 
\item 吻，Csaba; M ü ller，Thomas G.；Farkas-takács，anikó; Moór，Attila; Protopapa，Silvia; Parker，Alex H.; 等人 (2024年11月)。“JWST/MIRI发现的矮行星 (136472) Makemake突出的中红外过量表明正在进行活动”。天体物理学杂志快报。976(1): l9。arXiv:2410.22544。Bibcode:2024ApJ...976L...9K。doi:10.3847/2041-8213/ad8dcb。S2CID273695912。 
\item 阿尔瓦雷斯-坎达尔，阿尔瓦罗; 索萨-费利西亚诺，安娜卡罗莱纳州; 马丁斯-菲略，沃尔特; 皮尼利亚-阿隆索，诺埃米; 奥尔蒂斯，何塞·路易斯 (2020年10月)。“X-shooter看到的矮行星Makemake”。皇家天文学会月刊。497(4):5473-5479。arXiv:2008.01472。Bibcode:2020MNRAS.497.5473A。doi:10.1093/mnras/staa2329。S2CID220961487。 
\item "AstDys (136472) Makemake星历表"。意大利比萨大学数学系。已存档从原来的2019年9月25日。已检索4月9日，2019。
\item 奥尔蒂斯，J.L.;西卡迪，B.；布拉加-里巴斯，F.；阿尔瓦雷斯-坎达尔，A.；Lellouch, E.;Duffard, R.;皮尼拉-阿隆索，N.；伊万诺夫，V。D.;Littlefair, S.P.;卡马戈，J。I.B.;阿萨芬，M.；Unda-Sanzana, E.;Jehin, E.;莫拉莱斯，N.；Tancredi, G.;吉尔-赫顿，R.；De La Cueva, I.;科尔克，J。P.;达席尔瓦·内托，D.N.;曼弗里德，J.；Thirouin, A.;古铁雷斯，P。J.;Lecacheux, J.;吉伦，M.；莫里，A.；可乐，F.；利桑德罗，J.；穆勒，T.；雅克，C.；韦弗，D.(2012)。“来自恒星掩星的矮行星Makemake的反照率和大气约束”。自然。491(7425):566-569。Bibcode:2012 Natur.491..566O。doi:10.1038/nature11597。高密度脂蛋白:2268/142198。PMID23172214。S2CID4350486。已存档从原来的2023年2月8日。已检索7月1日，2019。  
\item 席林，总督(2009)。寻找行星X: 新世界和冥王星的命运(PDF)。斯普林格.doi:10.1007/978-0-387-77805-1。ISBN  978-0-387-77805-1。已检索10月26日，2025。
\item 施瓦姆，梅根。“探索相机”。La Silla QUEST KBO调查。耶鲁大学。已存档从原来的2015年7月22日。已检索10月23日，2025。
\item 布朗，迈克尔E。“发现厄里斯，已知最大的矮行星”。web.gps.caltech.edu。加州理工学院。已存档从2012年5月17日起。已检索10月17日，2025。
\item 布朗，M.E.;特鲁希略，C。A.;拉比诺维茨，D。(2005年7月29日)。格林，丹尼尔·W.E、 (编)。“2003 EL_61、2003 UB_313和2005 FY_9”。天文学联盟通告。8577(8577)。中央天文电报局: 1.Bibcode:2005IAUC.8577....1B。已检索10月23日，2025。
\item 布朗，迈克 (2010年12月)。我如何杀死冥王星，为什么它会来。Spiegel & Grau。ISBN 978-0-385-53108-5。
\item 布朗，迈克 (2013年4月3日)。“没有得到尊重的矮行星”。迈克·布朗的行星。存档自原来的2019年10月27日。已检索10月26日，2025。
\item 赫克特，杰夫 (2005年7月30日)。“在外太阳系发现的第十颗行星”。新科学家。存档自原来的2023年3月6日。已检索10月26日，2025。
\item 布朗，M.E.;特鲁希略，C。A.;拉比诺维茨，D.；马斯登，B。G.(2005年7月29日)。"MPEC 2005-O42: 2005 FY9"。小行星电子圆。2005年-o42。小行星中心。Bibcode:2005MPEC....O...42B。已检索10月23日，2025。
\item 巴内特，阿曼达 (2025年4月24日)。"Makemake事实"。NASA。已存档从2025年5月7日开始。已检索9月15日，2025。
\item “Makemake-或复活节兔子-进入太空名称之书”。每日空间。2008年7月15日已存档从原来的2008年8月1日。已检索10月31日，2025。
\item 辛诺特，罗杰W.(2005年9月11日)。“冥王星得到了一个小行星号”。天空和望远镜。已存档从原来的2020年8月8日。已检索10月29日，2025。
\item 比，Ker (2005年9月11日)。“冥王星现在只是一个数字: 134340”。Space.com。已存档从2011年5月23日起。已检索10月29日，2025。
\item Courtland，Rachel (2008年7月14日)。遥远的太阳系物体被命名为 “makemake”"。新科学家。存档自原来的2023年3月5日。已检索11月1日，2025。
\item 耶格尔，阿什利 (2008年7月21日)。“Makemake制作列表”。科学新闻。存档自原来的2022年5月20日。已检索10月30日，2025。
\item “命名天文物体-小行星”。国际天文学联合会。存档自原来的2013年7月7日。已检索10月31日，2025。远远超出海王星轨道的物体，包括矮行星，预计将被赋予与创造有关的神灵或人物的名称; 例如Makemake，人类的波利尼西亚创造者和生育之神...”“足够在海王星轨道之外的物体，可以在太阳系寿命的很大一部分时间内合理地确保轨道稳定性 (所谓的Cubewanos或“ 经典 ”TNOs) 被赋予与创造相关的神话名称。
\item "建议的新字符: 管道"。存档自原来的2022年1月29日。已检索1月29日，2022年。
\item Boeck，Moore (2018年1月30日)。“太阳系符号”。NASA.已存档从原来的2023年10月10日。已检索10月31日，2025。
\item 米勒，柯克 (2021年10月26日)。"Unicode请求矮行星符号"(PDF)。unicode.org。已存档(PDF)从原来的2022年3月23日。已检索8月6日，2022年。
\item JPL/NASA (2015年4月22日)。“什么是矮行星？”。喷气推进实验室。已存档从原始的2021年1月19日。已检索9月24日，2021。
\item 安德森，黛博拉 (2022年5月4日)。走出这个世界: 新的天文学符号被批准为Unicode标准。unicode.org。Unicode联盟。已存档从原来的2022年8月6日。已检索8月6日，2022年。
\item “JPL地平线136472 Makemake的在线星历 (2005 FY9) 在epochs JD 2450000.5-2460000.5”。JPL Horizons在线星历系统。喷气推进实验室。已检索10月31日，2025。使用Sun的解决方案。星历类型: 元素和中心: @ sun)
\item Moltenbrey，迈克尔 (2015年10月)。“矮行星”。小世界的黎明: 矮行星，小行星，彗星。天文学家的宇宙斯普林格.pp. 175-214。doi:10.1007/978-3-319-23003-0_5。ISBN 978-3-319-23003-0。已检索11月2日，2025。
\item “JPL Horizons 136472 Makemake (2005 FY9) 的在线星历，从2025年11月1日到2033年12月1日”。JPL Horizons在线星历系统。喷气推进实验室。已检索11月1日，2025。(r = 距太阳的距离，以AU为表示; rdot = 距太阳的径向速度，以km/s为表示)
\item “JPL Horizons 136472 Makemake (2005 FY9) 从2033年1月1日到2033年12月1日的在线星历”。JPL Horizons在线星历系统。喷气推进实验室。已检索11月1日，2025。(当径向速度或 “rdot” 从正变为负时，就会发生远日点。“hecl-lat” 是日心黄道纬度,或黄道上方的角距离。)
\item 帕克，A.H、; Buie，M。W.;格伦迪，W。M、；诺尔，K。S。(2016年4月26日)。"MPEC 2016-H46: S/2015 (136472) 1"。小行星电子圆。2016-h46。小行星中心。Bibcode:2016MPEC....H...46P。已检索9月19日，2025。
\item “JPL地平线136472 Makemake (2005 FY9) 的在线星历，从2103-jan-01到2104-dec-01”。JPL Horizons在线星历系统。喷气推进实验室。已检索11月1日，2025。(黄道穿越发生在日心黄道纬度或 “hecl-lat” 从正变为负。)
\item Muñoz-Gutiérrez, Marco A.;安东尼·佩姆伯特; 马修·J·莱纳；王尚宇 (2021年10月)。“外太阳系的长期动力稳定性。I.34个最大的跨海王星天体的规则和混沌演化“。天文杂志。162(4): 164。arXiv:2107.00240。Bibcode:2021AJ....162..164M。doi:10.3847/1538-3881/ac1102。S2CID235694711。 
\item 巴鲁奇，M.A.;Dalle Ore, C.;Fornasier, S.(2021年7月)。“作为冥王星背景的跨海王星天体: 天文视角”。新视野之后的冥王星系统。亚利桑那大学出版社。pp. 21-52。doi:10.2458/azu_uapress_9780816540945-ch003。ISBN 978-0-8165-4210-9。
\item 梅根·E·施万布；迈克尔·E·布朗；弗雷泽，韦斯利·C。(2014年1月)。“来自OSSOS的热主柯伊伯带尺寸分布”。天文杂志。147(1): 2。arXiv:1310.7049。Bibcode:2014AJ....147....2S。doi:10.1088/0004-6256/147/1/2。S2CID28744257。 
\item 佩蒂特，让-马克; 格拉德曼，布雷特; 卡夫拉斯，J。J.;米歇尔·T·班尼斯特；亚历山大·迈克; 沃克，凯瑟琳; 陈英东 (2023年4月)。“来自OSSOS的热主柯伊伯带尺寸分布”。天体物理学杂志快报。947(1): l4。arXiv:2309.01530。Bibcode:2023ApJ...947L...4P。doi:10.3847/2041-8213/acc525。S2CID258078423。 
\item 马克·W.Buie(2008年4月5日)。“136472的轨道配合和天体测量记录”。SwRI (空间科学系)。已存档从原来的2020年5月27日。已检索7月13日，2008。
\item 布朗，M.E.;特鲁希略，C。A.;拉比诺维茨，D。L.(2005年12月)。“在分散的柯伊伯带发现一个行星大小的物体”。天体物理学杂志。635(1):L97-L100。arXiv:astro-ph/0508633。Bibcode:2005ApJ...635L..97B。doi:10.1086/499336。S2CID1761936。  
\item M。E.布朗; K.M。巴克松; G。L.布莱克; E.L.Schaller; 等人 (2007年)。“2005 FY明亮的柯伊伯带天体上的甲烷和乙烷9"。天文杂志。133(1):284-289。Bibcode:2007AJ....133..284B。doi:10.1086/509734。S2CID12146168。 
\item 威廉姆斯，大卫R.(2024年1月11日)。"冥王星概况介绍"。NASA空间科学数据协调档案。NASA.存档自原来的2025年2月4日。已检索9月16日，2025。
\item Koberlein，Brian (2025年2月18日)。“甚至Eris和Makemake也可能有地热活动”。今天的宇宙。已检索9月15日，2025。
\item 卡尼，斯蒂芬 (2025年9月3日)。“我们太阳系中的行星大小和位置”。NASA。已存档从原来的2025年9月30日。已检索9月15日，2025。
\item 吻，C.；马顿，G.；帕克，A.H、; 格伦迪，W.；Farkas-Takacs, A.;Stansberry，J.; 等人 (2019年12月)。“矮行星的质量和密度 (225088) 2007或10”。伊卡洛斯。334:3-10。arXiv:1903.05439。Bibcode:2019年Icar .. 334 .... 3k。doi:10.1016/j.icarus.2019.03.013。S2CID119370310。 
\item “2022年CODATA值: 牛顿引力常数”。关于常数、单位和不确定性的NIST参考。NIST。2024年5月。已检索5月18日，2024。
\item 帕克，亚历克斯·H (2021年7月)。“跨尼普特尼亚空间与后冥王星范式”。新视野之后的冥王星系统。亚利桑那大学出版社。第1页。arXiv:2107.02224。doi:10.2458/azu_uapress_9780816540945-ch023。ISBN 978-0-8165-4210-9。S2CID 229245689。
\item E.L.夏勒; M.E.布朗 (2007年4月10日)。“柯伊伯带物体上的挥发性损失和保留”。天体物理学杂志。659(1):L61-L64。Bibcode:2007ApJ...659L..61S。doi:10.1086/516709。S2CID10782167。  
\item M。E.布朗; E.L.夏勒; G.A.布莱克 (2015)。“矮行星Makemake上的辐照产品”(PDF)。天文杂志。149(3): 105。Bibcode:2015AJ....149..105B。doi:10.1088/0004-6256/149/3/105。S2CID39534359。存档自原来的(PDF)于2020年4月12日。  
\item 格伦迪，W。M、；Wong, I.;格莱因，C.R、；Protopapa, S.;霍勒，B。J.;厨师，J。C.; 等人 (2024年3月)。“D/H的测量和13C/12Eris和Makemake上甲烷冰的C比: 内部活动的证据 “。伊卡洛斯。411115923: 115933。arXiv:2309.05085。Bibcode:2024Icar .. 41115923克。doi:10.1016/j.icarus.2023.115923。S2CID261681920。 
\item 佩尔纳，D.；Hromakina, T.;梅林，F.；Ieva, S.;Fornasier, S.;Belskaya, I.;Epifani, E.Mazzotta (2017年4月21日)。“矮行星Makemake的非常均匀的表面”(PDF)。皇家天文学会月刊。466(3):3594-3599。Bibcode:2017MNRAS.466.3594P。doi:10.1093/mnras/stw3272。ISSN0035-8711。已存档(PDF)从原始的2021年8月31日。已检索7月1日，2019。  
\item 巴鲁奇，M.A.;布朗，M.E.;埃默里，J。P.;梅林，F.(2008)。“跨底天体和半人马的组成和表面特性”(PDF)。超越海王星的太阳系。亚利桑那大学出版社。pp. 143-160。Bibcode:2008ssbn.book .. 143B。ISBN  978-0-8165-2755-7。S2CID 62885604。已存档 (PDF)从2017年8月12日的原始。
\item 利桑德罗，J.；皮尼拉-阿隆索，N.；佩达尼，M.；奥利瓦，E.；Tozzi, G.P.;格伦迪，W。M。(2006年1月)。“大型TNO富含甲烷冰的表面2005财年9: 跨尼普特尼亚带的冥王星双胞胎？”。天文学与天体物理学。445(3):L35-L38。Bibcode:2006A & A...445L..35L。doi:10.1051/0004-6361:200500219。S2CID56343540。  
\item A.N.海因策; D.德拉洪塔 (2009)。“柯伊伯带矮行星136472 Makemake (2005 FY9) 的自转周期和光曲线幅度”。天文杂志。138(2):428-438。Bibcode:2009AJ....138..428H。doi:10.1088/0004-6256/138/2/428。
\item “冥王星在真实的颜色”。当天的天文图片。NASA.2024年1月28日。已检索9月20日，2025。
\item 埃默里，J。P.;Wong, I.;布鲁内托，R.；厨师，J。C.;皮尼拉-阿隆索，N.；斯坦斯伯里，J。A.; 等人 (2024年3月)。“三颗矮行星的故事: 来自JWST光谱学的塞德纳、贡贡和夸尔上的冰和有机物”。伊卡洛斯。414116017。arXiv:2309.15230。Bibcode:2024Icar .. 41416017E。doi:10.1016/j.icarus.2024.116017。S2CID263152655。 
\item 穆西斯，奥利维尔; 亚伦，韦伦; 贝内斯特·库兹努，汤姆; 施尼伯格，安托ine (2025年4月)。“D/H比支持的Eris和Makemake上甲烷的原始起源”。天体物理学杂志快报。983(1): l12。arXiv:2503.15652。Bibcode:2025ApJ...983L..12M。doi:10.3847/2041-8213/adc134。
\item 洛伦兹，V.；皮尼拉-阿隆索，N.；Licandro, J.(2015年5月1日)。“矮行星 (136472) Makemake的旋转分辨光谱”。天文学与天体物理学。577: a86。arXiv:1504.02350。Bibcode:2015A & A...577A..86L。doi:10.1051/0004-6361/201425575。ISSN0004-6361。S2CID119253105。  
\item 克里斯托弗·R·格莱因；威廉·M·格伦迪；鲁尼，乔纳森I.；Wong，Ian; Protopapa，Silvia; Pinilla-alonso，Noemi; 等人 (2024年4月)。“Eris和Makemake上甲烷冰的中等D/H比作为其内部热液或变质过程的证据: 地球化学分析”。伊卡洛斯。412115999。arXiv:2309.05549。Bibcode:2024Icar .. 41215999G。doi:10.1016/j.icarus.2024.115999。S2CID261696907。 
\item 麦金农，W.B.;普里亚尔尼克，D.；斯特恩，S.A.；科拉迪尼，A。(2008)。“柯伊伯带天体和矮行星的结构和演化”(PDF)。超越海王星的太阳系。亚利桑那大学出版社。pp. 213-241.Bibcode:2008年ssbn。书 .. 213米。ISBN  978-0-8165-2755-7。S2CID 59421806。存档自原来的 (PDF)2025年11月6日
\item 帕克，亚历克斯 (2016年5月2日)。“Makemake的月亮”。planetary.org。行星社会。存档自原来的2018年10月21日。已检索5月2日，2016。
\item 乔纳森·奥卡拉汉 (2024年11月5日)。遥远的矮行星Makemake可能有一个令人惊讶的冰火山。新科学家。存档自原来的2024年12月22日。已检索9月19日，2025。
\item 安德鲁斯，罗宾·乔治 (2024年2月28日)。JWST间谍在冰冻的太阳系世界中出现令人惊讶的温暖迹象。科学美国人。已存档从原来的2024年2月29日。已检索9月19日，2025。
\item 达蒙德本宁菲尔德 (2024年3月27日)。“矮行星显示出最近地质活动的证据”。Eos。第105卷。doi:10.1029/2024EO240143。已存档从原来的2024年4月1日。已检索9月19日，2025。
\item “哈勃发现月球环绕矮行星Makemake”。NASA.2016年4月26日。已检索9月19日，2025。
\item 托马斯，迈克W。(2025年9月9日)。“SwRI领导的团队在Makemake上发现了甲烷气体”。西南研究院。已存档从原来的2025年9月10日。已检索9月19日，2025。
\item 斯蒂芬妮·M·门滕；迈克尔·M·索里；布拉姆森，阿里M.(2022年8月)。“Charon和其他柯伊伯带物体上的内源挥发物”。自然通讯。13(1) 4457。Bibcode:2022年纳科。。13.4457米。doi:10.1038/s41467-022-31846-8。PMC9363412。PMID35945207。S2CID251469650。   
\item 卢克·D·伯克哈特；Ragozzine，达林; 布朗，迈克尔E。(2016年6月)。“在矮行星Haumea周围深入寻找其他卫星”。天文杂志。151(6): 162。arXiv:1605.03941。Bibcode:2016AJ....151..162B。doi:10.3847/0004-6256/151/6/162。S2CID23744219。 
\item 布鲁诺·西卡迪; 费利佩·布拉加·里巴斯; 马克·W·布伊；奥尔蒂斯，何塞·路易斯; 罗克斯，弗朗索瓦 (2024年12月)。“跨尼普顿天体的恒星掩星”。天文学和天体物理学评论。32(1) 6。arXiv:2411.07026。Bibcode:2024A和ARv。。32。6s。doi:10.1007/s00159-024-00156-x。S2CID273963101。6. 
\item Schmid，Deb (2025年2月15日)。SwRI科学家在冰冷的矮行星中发现了地热活动的证据。西南研究院。已存档从原来的2025年2月15日。已检索9月27日，2025。
\item 德索萨·里贝罗，拉斐尔; 莫比德利，亚历山德罗; 雷蒙德，肖恩·N；伊齐多罗，安德烈; 戈麦斯，罗德尼; 维埃拉·内托，埃内斯托 (2020年3月)。“早期巨型行星不稳定性的动力学证据”。伊卡洛斯。339113605。arXiv:1912.10879。Bibcode:2020年Icar .. 33913605D。doi:10.1016/j.icarus.2019.113605。S2CID209444878。 
\item 杰弗里·M·摩尔；威廉·B·麦金农.(2021年5月)。“地质学上多样化的冥王星和Charon: 对柯伊伯带矮行星的影响”。年度审查。49:173-200。Bibcode:2021年，49..173m。doi:10.1146/annurev-earth-071720-051448。S2CID229420916。 
\item "MPC数据库搜索: H < 3.8和a > 20 AU"。小行星中心。已存档从原来的2025年8月16日。已检索8月16日，2025。22个比H = 3.8更亮的海王星天体(H值越小越亮)
\item 艾伦·W·哈里斯；哈里斯，艾伦W.(1997年4月)。“关于辐射反照率和小行星直径的修订”。伊卡洛斯。126(2):450-454。Bibcode:1997年Icar。。126。450小时。doi:10.1006/ICAR.1996.5664。S2CID123488868。 
\item 福特，多米尼克 (2025年3月30日)。“136472反对派的Makemake”。In-The-Sky.org。已存档从原来的2025年4月1日。已检索10月17日，2025。
\item “JPL地平线在线星历136472 Makemake (2005 FY9) 于2005年1月1日至20 1月1日 (显示视星等和星座)”。JPL Horizons在线星历系统。喷气推进实验室。已检索10月17日，2025。星历类型: 观察者和中心: 500 @ 399 (地心)
\item 斯科特·S·谢泼德；Udalski，Andrzej; Trujillo，Chadwick; Pietrzynski，Marcin; Bellerose，Grzegorz; Poleski，Radoslaw; 等人 (2011年10月)。“明亮的柯伊伯带天体的南部天空和银河系平面调查”。天文杂志。142(4): 98。arXiv:1107.5309。Bibcode:2011AJ....142...98S。doi:10.1088/0004-6256/142/4/98。S2CID53552519。 
\item Tichy，Milos; Ticha，Jana (2007年1月17日)。"TNO 2005 FY9"。KLENOT项目。克莱特天文台。已存档从2011年6月15日起。已检索10月16日，2025。
\item 贝茨，布鲁斯 (2007年3月1日)。“Shoemaker NEO补助金的过去接受者的更新 (2007年3月1日)”。planetary.org。行星社会。已存档从2025年6月17日开始。已检索10月16日，2025。
\item "M.P.S.138898"(PDF)。小行星和彗星补充(138898)。小行星中心: 110。2005年8月7日。已检索10月16日，2025。
\item "M.P.S.598178"(PDF)。小行星和彗星补充(598178)。小行星中心: 494。2015年4月12日。已检索10月23日，2025。
\item “矮行星Makemake缺乏大气层”。欧洲南方天文台。已存档从原来的2017年1月18日。已检索11月23日，2012。
\item 奥尔蒂斯，J.L.;西卡迪，B.；阿萨芬，M.；阿尔瓦雷斯-坎达尔，A.；伊万诺夫，V。D.;Camargo，J.; 等人 (2011年10月)。2011年4月23日Makemake的恒星掩星(PDF)。2011年epsc-dps联席会议。法国南特。第704页。Bibcode:2011epsc.conf .. 704O。
\item 奥尔蒂斯，J.L.;西卡迪，B.；卡马戈，J。I.B.;桑托斯-桑兹，P.；布拉加-里巴斯，F.(2019年5月)。“跨内天体的恒星掩星: 从预测到观测和未来展望”。arXiv:1905.04335[astro-ph.EP]。
\item 格伦迪，W。M、；麦金农，W.B.;Ammannito, E.;Aung, M.;Bellerose，J; Brenker，F.; 等人 (2009年12月)。2013-2022年冰矮行星的探索策略(PDF)。小机构评估小组。月球和行星研究所。Bibcode:2009AGUFM.P43D1471G。
\item 麦格拉纳汉，R.；萨根，B.；鸽子，G.；图洛斯，A.；Lyne, J.E.;埃默里，J。P。(2011)。“对跨Neptunian物体的任务机会的调查”。英国星际学会杂志。64:296-303。Bibcode:2011JBIS...64..296M。已存档(PDF)从2025年10月19日开始。
\item 迦南的安德森; 米尔斯的约翰逊；诺亚·C·纳什；五旬节，米切尔；约瑟夫·A·锡拉丘兹；冯·巴根，劳拉·E；詹姆斯·E·莱恩.(2024年8月)。动力木星的重力协助缩短了跨Neptunian物体的运输时间。2024年AAS/AIAA天体动力学专家会议。丹佛，科罗拉多州。p. AAS 24-472。已检索10月19日，2025。
\item 阿曼达·M·赞加里；蒂芙尼·J·芬利；斯特恩，S.艾伦; 塔普利，马克·B。(2019年5月)。“重返柯伊伯带: 2025年至2040年的发射机会”。宇宙飞船和火箭杂志。56(3):919-930。arXiv:1810.07811。Bibcode:2019年JSpRo .. 56 .. 919Z。doi:10.2514/1.A34329。S2CID119033012。 
\item 巴鲁奇，M.A.;Fulchignoni, M.(1982年8月)。“小行星光曲线对小行星方向参数和形状的依赖性”。地球、月球和行星。27: 47。Bibcode:1982年M & P。27...47B。doi:10.1007/BF00941556。
\end{enumerate}
\subsection{外部链接}
\begin{itemize}
\item 冰冷矮行星Eris和Makemake上的地热活动-行星收音机 播客主持人:行星社会(2024年3月6日)
\item 2012年Kavli奖得主演讲: 来自外太阳系的故事迈克尔·E。棕色 (视频)
\item Makemake天空图和坐标-天空直播
\item Precovery图像用1.06米Kle ť 天文台2003年4月20日望远镜
\item Makemake图像由Keck 1望远镜2010年2月18日 (Mike Brown)
\item Makemake在AstDyS-2，小行星-动态站点
\item 星历 · 观测预测 · 轨道信息 · 适当的元素 · 观测信息
\item Makemake在JPL小型数据库 在Wikidata上编辑此内容
\item 关闭进近 · 发现 · 星历 · 轨道查看器 · 轨道参数 · 物理参数
\end{itemize}